\section{Gauge, Rank Strata, and Moving Information}
\label{sec:stratified-geometry}

The previous section fixes \(A\).  We now let the information channel move.
Relative to \(P_0\), the pair \((A,R)\) has a change-of-visible-coordinates
gauge
\[
  (A,R)\cdot B=(AB,B^\top RB),
  \qquad B\in GL(m).
\]
First define the total, generally redundant descriptor space
\begin{equation}
  \widetilde{\mathfrak D}_m(P_0)
  =
  \left\{(A,R):\rank A=m,\ R\in\SPD^m\right\}/GL(m),
  \label{eq:total-descriptor-space}
\end{equation}
and the everywhere-defined completion parameterization
\begin{equation}
  \widetilde\Gamma_{m,P_0}([A,R])=\Psi_{P_0,A}(R).
  \label{eq:total-completion-map}
\end{equation}
When \(m<d\), it is redundant if some visible modes equal their reference
values, because those modes can rotate without changing the completion.  For
\(m=d\), the Grassmann factor is a point and this rotational redundancy is
absent.  The regular subset is
\begin{equation}
  \mathfrak D_m^{\mathrm{reg}}(P_0)
  =
  \left\{(A,R):\rank A=m,\ R\in\SPD^m,
  \ \det(R-A^\top P_0A)\ne0\right\}/GL(m),
  \label{eq:descriptor-space}
\end{equation}
and the rank stratum
\begin{equation}
  \mathfrak S_m(P_0)
  =\{P\in\SPD^d:\rank(P-P_0)=m\}.
  \label{eq:rank-stratum}
\end{equation}

\begin{theorem}[Completion bundle and complete rank stratification]
\label{thm:rank-stratification}
The map
\begin{equation}
  \Gamma_{m,P_0}([A,R])=\Psi_{P_0,A}(R)
  \label{eq:stratum-map}
\end{equation}
is a diffeomorphism from
\(\mathfrak D_m^{\mathrm{reg}}(P_0)\) onto
\(\mathfrak S_m(P_0)\).  Its inverse descriptor is intrinsic:
\begin{equation}
  \range(A)=P_0^{-1}\range(P-P_0),
  \qquad
  R=A^\top PA,
  \label{eq:intrinsic-descriptor}
\end{equation}
up to the stated \(GL(m)\) action.  The strata exhaust and close according to
\begin{equation}
  \SPD^d=\bigsqcup_{m=0}^d\mathfrak S_m(P_0),
  \qquad
  \overline{\mathfrak S_m(P_0)}^{\,\SPD^d}
  =\bigcup_{r=0}^m\mathfrak S_r(P_0).
  \label{eq:strata-closure}
\end{equation}
\end{theorem}

The pullback AIRM seminorm on the total descriptor space resolves the cost of
changing the visible target and the visible subspace.  It becomes a genuine
quotient metric on the regular subset, and on the entire descriptor space when
\(m=d\).  Normalize a representative by
\(A^\top P_0A=I_m\), use
the horizontal gauge \(A^\top P_0\dot A=0\), choose \(N\) so that
\([P_0^{1/2}A,N]\) is orthogonal, and put
\(K=N^\top P_0^{1/2}\dot A\).

\begin{theorem}[Closed-form completion-pair metric and regular quotient geometry]
\label{thm:stratified-line-element}
At any normalized representative of
\(\widetilde{\mathfrak D}_m(P_0)\), represent two horizontal quotient
directions by
\[
  \xi_i=(H_i,K_i),
  \qquad H_i\in\Sym^m,
  \qquad K_i=N^\top P_0^{1/2}\dot A_i,
  \qquad i\in\{1,2\}.
\]
Then the pullback AIRM pair kernel is
\begin{equation}
  \boxed{
  \begin{aligned}
  \mathfrak G_R^{\mathrm{comp}}(\xi_1,\xi_2)
  &:={}
  \left\langle
  D\widetilde\Gamma_{m,P_0}[\xi_1],
  D\widetilde\Gamma_{m,P_0}[\xi_2]
  \right\rangle_{\widetilde\Gamma_{m,P_0},\AI}
  \\
  &={}
  \tr(R^{-1}H_1R^{-1}H_2)
  +2\tr\!\left(K_1\Xiop(R)K_2^\top\right),
  \end{aligned}}
  \label{eq:completion-pair-kernel}
\end{equation}
where
\begin{equation}
  \Xiop(R)=(R-I)R^{-1}(R-I)=R+R^{-1}-2I\succeq0.
  \label{eq:mismatch-weight}
\end{equation}
In particular, for \(\xi=(\dot R,K)\),
\begin{equation}
  \boxed{
  \norm{D\widetilde\Gamma_{m,P_0}[\dot A,\dot R]}_{\widetilde\Gamma_{m,P_0},\AI}^2
  =\mathfrak G_R^{\mathrm{comp}}(\xi,\xi)
  =
  \fro{R^{-1/2}\dot R R^{-1/2}}^2
  +2\tr\!\left(K\Xiop(R)K^\top\right),}
  \label{eq:stratified-line-element}
\end{equation}
so visible-target and subspace-motion blocks are orthogonal.

More explicitly, define the completion feature
\begin{equation}
  \Phi_R(\xi)
  =
  \left(
  R^{-1/2}HR^{-1/2},
  \sqrt{2}\,K(R-I)R^{-1/2}
  \right).
  \label{eq:completion-feature}
\end{equation}
Then
\(
\mathfrak G_R^{\mathrm{comp}}(\xi_i,\xi_j)
=\langle\Phi_R(\xi_i),\Phi_R(\xi_j)\rangle_F
\)
in the direct-sum Frobenius space.  Hence for any directions
\(\xi_1,\ldots,\xi_k\), the pair matrix
\begin{equation}
  \mathbf G
  =\left[\mathfrak G_R^{\mathrm{comp}}(\xi_i,\xi_j)\right]_{i,j=1}^k
  \succeq0,
  \qquad
  \rank\mathbf G
  =\dim\operatorname{span}\{\Phi_R(\xi_i):1\le i\le k\}.
  \label{eq:completion-gram-rank}
\end{equation}
Moreover,
\begin{equation}
  \mathfrak G_R^{\mathrm{comp}}(\xi_i,\xi_j)=0
  \quad\Longleftrightarrow\quad
  \Phi_R(\xi_i)\perp\Phi_R(\xi_j).
  \label{eq:completion-pair-orthogonality}
\end{equation}
The exact horizontal differential kernel of the total parameterization is
\begin{equation}
  \ker D\widetilde\Gamma_{m,P_0}
  =\{(\dot A,0):K(R-I)=0\}.
  \label{eq:stratified-kernel}
\end{equation}
If \(m<d\), the pullback is nondegenerate precisely on the regular subset
\(\det(R-I)\ne0\), where
\(\widetilde\Gamma_{m,P_0}=\Gamma_{m,P_0}\) is the diffeomorphism of
\cref{thm:rank-stratification}.  At singular descriptors, rotation of a
reference-valued visible mode is unidentifiable from the completed matrix.
If \(m=d\), there is no subspace-rotation block: the visible AIRM term is
nondegenerate for every \(R\), and the total completion map is a
diffeomorphism onto \(\SPD^d\).  Thus descriptor singularity occurs only when
\(m<d\); the ambient SPD manifold and its AIRM remain smooth and
nondegenerate in every case.
If
\(\operatorname{dist}(\operatorname{spec}(R),1)\ge\zeta>0\) and
\(0<a\le\lambda(R)\le b\), then subspace motion is quantitatively controlled:
\begin{align}
  \fro{R^{-1/2}\dot R R^{-1/2}}^2
  +\frac{2\zeta^2}{b}\fro K^2
  &\le \norm{\dot P}_{P,\AI}^2
  \nonumber\\
  &\le
  \fro{R^{-1/2}\dot R R^{-1/2}}^2
  +2\max_{x\in[a,b]}\frac{(x-1)^2}{x}\fro K^2.
  \label{eq:stratified-conditioning}
\end{align}
\end{theorem}

\begin{corollary}[Reduced completion value and completed motion]
\label{cor:reduced-completion-value}
At fixed \(P_0\), in the normalized chart \(A^\top P_0A=I\), the reduced
minimum-distance value is
\begin{equation}
  \overline E(A,R)
  :=\min_{A^\top PA=R}\frac12\dai(P_0,P)^2
  =\frac12\dai(I,R)^2.
  \label{eq:reduced-completion-energy}
\end{equation}
It contains no independent information-subspace coordinate.  Consequently,
holding \(R\) fixed, all first and mixed parameter derivatives generated by
horizontal variations of \(A\) vanish.  Exact completion has already removed
the need for an implicit vertical stationarity solve: the minimizing ambient
matrix is the explicit section \(P=\Psi_{P_0,A}(R)\).  The pullback pair
metric \eqref{eq:completion-pair-kernel} is a different second-order object:
it measures motion of this minimizing matrix rather than curvature of the
reduced value.
\end{corollary}

\begin{remark}[Reduced value versus completion motion]
\label{rem:value-versus-motion}
The distinction is already visible for \(d=2\), \(m=1\), \(P_0=I\), and a
scalar target \(r>0\), \(r\ne1\).  With
\(u(\theta)=(\cos\theta,\sin\theta)^\top\), the canonical completion is
\[
  P_\star(\theta)=I+(r-1)u(\theta)u(\theta)^\top.
\]
Its reduced minimum-distance value is the constant
\(\tfrac12\log^2r\), whereas
\[
  \norm{\dot P_\star(0)}_{P_\star(0),\AI}^2
  =\frac{2(r-1)^2}{r}>0.
\]
Thus vanishing subspace derivatives of the optimized value coexist with
nonzero metric motion of its canonical minimizer.  A frozen--relaxed
second-variation decomposition may likewise contain nonzero direct and
Schur terms that cancel; their split depends on the chosen trivialization of
the moving constraint.  The completion kernel itself is analytic for every
\(R\in\SPD^m\), including repeated eigenvalues, and its degeneracy is exactly
the descriptor non-identifiability in \eqref{eq:stratified-kernel}.
\end{remark}

The same identity has a gauge-covariant moving-reference form.  Let
\(P_{0,t}=G_tG_t^\top\), define
\[
  R_{0,t}=A_t^\top P_{0,t}A_t,
  \quad
  U_t=G_t^\top A_tR_{0,t}^{-1/2},
  \quad
  \bar R_t=R_{0,t}^{-1/2}R_tR_{0,t}^{-1/2},
\]
\[
  Q_t=I+U_t(\bar R_t-I)U_t^\top,
  \qquad
  \Omega_t=G_t^{-1}\dot G_t,
  \qquad
  \mathfrak D_tQ_t=\dot Q_t+\Omega_tQ_t+Q_t\Omega_t^\top.
\]

\begin{proposition}[Covariant moving-reference identity]
\label{prop:moving-reference}
\mbox{}\par\noindent
For \(P_t=\Psi_{P_{0,t},A_t}(R_t)=G_tQ_tG_t^\top\),
\begin{equation}
  \boxed{
  \norm{\dot P_t}_{P_t,\AI}^2
  =\fro{Q_t^{-1/2}(\mathfrak D_tQ_t)Q_t^{-1/2}}^2.}
  \label{eq:moving-reference}
\end{equation}
Under \(G_t'=G_tO_t\) with \(O_t\) orthogonal,
\[
  Q_t'=O_t^\top Q_tO_t,
  \qquad
  \mathfrak D_t'Q_t'=O_t^\top(\mathfrak D_tQ_t)O_t.
\]
Hence \eqref{eq:moving-reference} is independent of the moving frame.
\end{proposition}

For a fixed reference and a fixed-rank curve, integrating
\eqref{eq:stratified-line-element} gives the exact energy decomposition
\begin{align}
  \int_0^T\norm{\dot P_t}_{P_t,\AI}^2\,\dd t
  &=\int_0^T\fro{R_t^{-1/2}\dot R_tR_t^{-1/2}}^2\,\dd t
  \nonumber\\
  &\quad+2\int_0^T\tr(K_t\Xiop(R_t)K_t^\top)\,\dd t.
  \label{eq:integrated-energy}
\end{align}

\begin{corollary}[Boundary of continuous rank loss]
\label{cor:rank-loss}
Suppose
\(P_t=P_0^{1/2}[I+U_t(R_t-I)U_t^\top]P_0^{1/2}\) has active rank \(m\)
for \(t<t_0\) and converges to a point of rank at most \(m-r\) relative to
\(P_0\).  Then at least \(r\) singular values of \(R_t-I\) converge to zero.
Thus continuous rank loss can occur only through visible modes approaching the
reference geometry.
\end{corollary}

The proofs are collected in \cref{app:stratified-proofs}.  The line element is
a kinematic identity on the completion bundle.  A learning or control law
would require an additional evolution rule for \((P_0,A,R)\).
